% Generated draft: requires mathematical and source review.
\documentclass[11pt]{article}
\usepackage[margin=25mm]{geometry}
\usepackage{fontspec}
\setmainfont{Latin Modern Roman}
\usepackage{amsmath,amssymb,mathtools,unicode-math}
\setmathfont{Latin Modern Math}
\usepackage{enumitem,needspace,xurl,hyperref}
\setlength{\parindent}{0pt}
\setlength{\parskip}{5pt}
\setlength{\emergencystretch}{3em}
\newcommand{\N}{\mathbb{N}}
\newcommand{\Z}{\mathbb{Z}}
\newcommand{\Q}{\mathbb{Q}}
\newcommand{\R}{\mathbb{R}}
\newcommand{\C}{\mathbb{C}}
\newcommand{\F}{\mathbb{F}}
\newcommand{\A}{\mathbb{A}}
\newcommand{\PP}{\mathbb{P}}
\newcommand{\E}{\mathbb{E}}
\newcommand{\eps}{\varepsilon}
\newcommand{\dd}{\,\mathrm{d}}
\begin{document}
\section*{UnsolvedMath ID: 11000254}
\textbf{Source problem code:} catalogue entry
\subsection{Definitions and mathematical statement}
For integers \(g,h\ge2\), let \(\Sigma_g,\Sigma_h\) be the source's genus-\(g\) and genus-\(h\) surfaces, conventionally closed and oriented. Question 4.7 asks for a short exact sequence
\[
 1\longrightarrow\pi_1(\Sigma_g)\longrightarrow\Gamma
       \longrightarrow\pi_1(\Sigma_h)\longrightarrow1
\]
with \(\Gamma\) word hyperbolic. Exactness means the first map is injective, the last is surjective, and its kernel is precisely the first image. Word hyperbolicity refers to uniformly thin geodesic triangles in a finite-generator Cayley graph. For fixed \(g,h\) these are the requested groups; the source's wording should be checked for whether it asks for every pair or existence for some pair in that range.

The extension gives an outer action of the quotient surface group on the kernel surface group. The source, by an argument from its Theorem 4.2, says this produces a purely pseudo-Anosov surface subgroup of \(\Gamma_g\). The notation \(\Gamma_g\), orientation conventions and theorem hypotheses are not supplied, so that consequence is retained as a source report rather than independently proved or used to alter the exact-sequence request.

The first contextual comment states that such a \(\Gamma\) cannot be the fundamental group of a closed complex-hyperbolic surface. It explains, through [20], [25] and [17], that an assumed realization \(X=\mathbb H^2_{\C}/\Gamma\) would yield a nonsingular holomorphic fibration \(X\to\Sigma_h\) inducing this exact sequence. The obstruction and precise hypotheses are not included. The second comment attributes a construction idea to Ian Agol and introduces a Thurston theorem in [40], but the theorem statement is absent and is not reconstructed.

The supplied summary links the problem to convex-cocompact monodromy and anticipates a positive construction, while reporting no established example of the requested word-hyperbolic surface-by-surface group. That expectation is not a theorem or a completed construction.
\subsection{Short English statement}
Can an extension of one closed genus-at-least-two surface group by another be word hyperbolic? The source connects this to purely pseudo-Anosov surface subgroups and reports a complex-hyperbolic obstruction, but the referenced proofs and Agol/Thurston construction idea are incomplete in the extract.
\subsection{Sources}
\begin{itemize}
\item Local UnsolvedMath registry, permanent ID \(11000254\), exact source code \texttt{catalogue entry}, Question 4.7, PDF page 289. No individual proposer or proposed year is supplied; the title is truncated.
\item Farb, edited volume \emph{Problems on Mapping Class Groups and Related Topics}: \url{https://www.math.uchicago.edu/~farb/papers/mcgbook.pdf}. The volume has 329 explicitly typeset numbered items across chapters, repeated local numbering distinguished by PDF page. The extract includes chapter/page furniture ``Surface subgroups of mapping class groups,'' page 283.
\item Source cross-reference Theorem 4.2; [20], [25] and [17] for the complex-hyperbolic/fibration discussion; Thurston [40] for a following theorem, introduced through an idea attributed to Ian Agol. Full entries and the latter theorem are absent.
\item Extraction pdf-font-structure, accessed 2026-07-29. Source evidence NEEDS\_REVIEW presents an open problem and calls for a current primary-literature audit. Supplied OPEN-TRIAGE summary reports the existence question open and discusses anticipated convex-cocompact monodromy without a verified construction or complete reference.
\item Dataset attribution: Ulam AI / UnsolvedMath Contributors, CC BY 4.0. Source rights NEEDS\_REVIEW: publicly accessible author/publisher copy, redistribution terms requiring review.
\end{itemize}
\end{document}
